\bookchapter{Глава 45: Гарантия качества — делаем код от ИИ надёжным}

\chapterintro{У кода, написанного ИИ, есть особенность: он обычно «выглядит правильно» — переменные названы по канону, комментарии на месте, логика прозрачна. Но «выглядит правильно» не значит «правильно на деле». Скрытые дефекты в коде ИИ часто заметить труднее, чем в написанном человеком: внешнее качество так высоко, что бдительность дремлет. В этой главе мы учимся выстраивать для кода ИИ действенную систему контроля качества — три линии обороны (тесты, ревью, контроль версий), которые превращают «выглядит правильно» в «правильно на деле».}

\booksection{Тестовая дисциплина: поручите ИИ писать тесты}

Тесты — первая линия обороны качества в программировании с ИИ. Хорошая новость: попросить ИИ написать тесты проще и надёжнее, чем заставить его писать функциональный код, — ожидаемое поведение теста известно заранее, и случай «выглядит правильно, а работает неверно» тут случается куда реже.

Процесс тест-ориентированной разработки с ИИ: сначала ИИ пишет функциональный код; как только базовая функциональность проверена, немедленно поручаем ему написать тест-кейсы. Тесты должны покрывать нормальный путь, граничные случаи и исключительные ситуации. Затем ИИ правит функциональный код по отзывам тестов. Плюс процесса: однажды написанные тесты становятся постоянным инструментом проверки — любое последующее изменение можно проверить запуском тестов и убедиться, что регрессий нет.

\begin{lstlisting}
Напишите модульные тесты для функции calculate_discount.
Тесты должны покрывать:
1. Нормальный путь:
   - товар за 100 юаней, скидка 20% -> 80
   - товар за 200 юаней, скидка 50% -> 100
2. Граничные случаи:
   - товар за 0 юаней, скидка 20% -> 0
   - товар за 100 юаней, скидка 0% -> 100 (без скидки)
   - товар за 100 юаней, скидка 10% -> 100 (полная цена)
3. Ошибочные случаи:
   - отрицательная цена -> должна выбрасываться ошибка
   - скидка > 1 -> должна выбрасываться ошибка
   - нечисловой ввод -> должна выбрасываться ошибка
Используйте фреймворк Vitest и дайте полный файл тестов.
\end{lstlisting}

Обратите внимание на особенность такой инструкции: она явно перечисляет три класса тест-кейсов (нормальные, граничные, исключения) и не оставляет ИИ простора для «выборочного тестирования». ИИ порой пишет тесты только нормального пути, чтобы «быстренько отчитаться»; явное перечисление всех классов позволяет этого избежать.

\booksection{Приёмы ревью кода}

Даже с покрытием тестами код, созданный ИИ, всё равно нужно просматривать самому. Ревью — это не построчное чтение: оно слишком долгое и слишком зависит от вашего уровня программирования. Действенное ревью кода ИИ должно смотреть на несколько ключевых измерений:

\begin{center}
\begin{tabularx}{\linewidth}{|X|X|X|}
\hline
\textbf{Измерение ревью} & \textbf{На что смотреть} & \textbf{Типичные проблемы} \tabularnewline
\hline
Безопасность & Валидация ввода, обращение с ключами, защита от SQL-инъекций & ИИ часто забывает валидировать ввод и зашивает ключи прямо в код \tabularnewline
Обработка ошибок & Перехватываются ли исключения, осмысленны ли сообщения об ошибках & ИИ любит пустые блоки catch или один console.log \tabularnewline
Производительность & Нет ли явных проблем с производительностью: N+1-запросы, бесконечные циклы & ИИ иногда выдаёт алгоритмы сложности O(n²) \tabularnewline
Сопровождаемость & Ясна ли структура кода, осмысленны ли имена & ИИ иногда генерирует чрезмерно абстрактный код \tabularnewline
Управление зависимостями & Нужны ли подключённые зависимости, зафиксированы ли версии & ИИ склонен тянуть лишние библиотеки \tabularnewline
\hline
\end{tabularx}
\end{center}

\begin{center}\small\color{gray}Ключевые измерения ревью кода ИИ\end{center}

Не обязательно проверять всё самому — первую проверку можно поручить ИИ. Инструкция вроде: «Проверь этот код с точки зрения безопасности, обработки ошибок и производительности, укажи потенциальные проблемы». Способность ИИ к самопроверке постоянно растёт; полностью заменить ревью человеком она не может, но отсеет большинство очевидных проблем.

\booksection{Контроль версий: роль Git в программировании с ИИ}

Роль Git в программировании с ИИ ещё важнее, чем в традиционном. Каждое вмешательство ИИ может принести непредсказуемые перемены, поэтому нужна надёжная «таблетка от раскаяния» — и контроль версий и есть эта таблетка.

Лучшая практика Git в работе с ИИ — «коммить часто, шагать маленькими шагами». Завершили проверяемый фрагмент функциональности — делайте коммит. Тогда, если очередная правка ИИ что-то сломает, вы сможете точно откатиться к последнему рабочему состоянию, а не перечёркивать всю работу.

Ещё один полезный приём — «изоляция в ветке». Для изменений ИИ с непредсказуемым результатом (например, крупного рефакторинга) сначала заведите отдельную ветку. Результат устроил — вливайте в основную ветку; не устроил — просто удаляйте её, не задев основную. Это куда безопаснее, чем много раз править основную ветку.

\begin{enumerate}
  \item Перед каждой правкой ИИ следите, чтобы рабочая копия была чистой (без незакоммиченных изменений). Тогда, если результат не понравится, можно быстро откатиться командой git checkout .
  \item После каждой правки ИИ, как только проверка пройдена, сразу делайте коммит. В сообщении опишите, что именно изменилось, например «feat: добавлена регистрация пользователей».
  \item Для крупного рефакторинга сначала создайте ветку: git checkout -b refactor/auth-system. Правьте в ветке и вливайте только после прохождения тестов.
  \item Возьмите в привычку смотреть diff. После каждой правки ИИ просматривайте командой git diff, что именно изменилось. Это не только помогает находить проблемы, но и учит вас коду.
\end{enumerate}

\booksection{Частые ловушки и как их обходить}

В программировании с ИИ есть несколько ловушек, в которые попадают почти все новички. Знать о них заранее — значит сберечь много времени и сил.

\booksubsection{Ловушка первая: ловушка доверия}

«Код от ИИ выглядит профессионально — значит, точно в порядке?» — самая опасная ловушка. Сгенерированный ИИ код снаружи очень опрятен — имена по канону, комментарии на месте, структура прозрачна, — и это «ощущение внешнего профессионализма» усыпляет бдительность. Но внешняя опрятность не гарантирует корректной логики. Противоядие: проверять всегда тестами, а не «по виду».

\booksubsection{Ловушка вторая: ловушка раздувания}

«Раз ИИ пишет за меня код, не сделать ли функции побольше?» — прямой путь к раздуванию кода. ИИ любит «дать сверх заказа»: попросите страницу входа — получите целую систему управления пользователями. Лишние функции не только усложняют проект, но и повышают шанс багов. Противоядие: строго следовать принципу Simplicity First и прямо говорить ИИ «делай только эту функцию, ничего не добавляй».

\booksubsection{Ловушка третья: ловушка понимания}

«ИИ ведь понимает, что я имею в виду?» — не факт. Способности ИИ к пониманию велики, но с расплывчатыми требованиями он всё же «догадывается». Догадается неверно — код будет неверным. Хуже того, ошибочный код обычно «выглядит вполне разумно», и кажется, что вас поняли. Противоядие: давать в инструкции достаточно конкретных деталей, оставляя ИИ меньше простора для догадок.

\booksubsection{Ловушка четвёртая: ловушка контекста}

«Я же говорил об этом в прошлой беседе — ИИ должен помнить?» — скорее всего, не помнит. С ростом длины беседы или открытием новой сессии ИИ теряет ранний контекст. Если ключевое ограничение не внесено в файл памяти проекта, при дальнейших правках его, скорее всего, проигнорируют. Противоядие: записывайте ключевые ограничения в файл памяти проекта, не полагаясь на память беседы.

\begin{tcbwarn}{Четыре ловушки одной строкой}
Ловушка доверия — не верьте «выглядит правильно»; ловушка раздувания — ИИ любит «дать сверх заказа»; ловушка понимания — с расплывчатыми требованиями он «догадывается»; ловушка контекста — память беседы «забывает». Противоядия соответственно: проверка тестами, отсечение лишнего, достаточно деталей, запись в файл памяти.

\end{tcbwarn}

\begin{tcbpractice}{Попробуйте сами: выстройте для своего проекта три линии обороны}
Если вы уже начали проект 1 (сайт инструментов): 1) попросите ИИ написать три класса тестов для ключевых функций; 2) просмотрите код по пяти измерениям (безопасность / обработка ошибок / производительность / сопровождаемость / зависимости); 3) инициализируйте проект в git и сделайте первый коммит. Когда три линии обороны готовы, проект по-настоящему становится вашим.

\end{tcbpractice}

\begin{tcbsummary}{Краткие итоги главы}
\begin{itemize}
  \item Код ИИ «выглядит правильно» $\neq$ «правильно на деле» — слишком высокое внешнее качество усыпляет бдительность
  \item Поручайте ИИ писать тесты: явно перечисляйте три класса кейсов (нормальные / граничные / исключения), чтобы не было «выборочного тестирования»
  \item Пять измерений ревью: безопасность, обработка ошибок, производительность, сопровождаемость, управление зависимостями
  \item Git — «таблетка от раскаяния»: коммитьте часто, шагайте маленькими шагами, изолируйтесь в ветках, почаще смотрите diff
  \item Четыре ловушки: доверие (проверка тестами), раздувание (отсечение лишнего), понимание (достаточно деталей), контекст (запись в файл памяти)
\end{itemize}

\end{tcbsummary}
